\documentclass[../../../include/open-logic-section]{subfiles}

\begin{document}

\olfileid[id]{sth}{cardinals}{hp}
\olsection{Lampiran: Prinsip Hume}

Dalam \olref[cp]{sec}, kita menguraikan Prinsip Cantor. Prinsip tersebut
berbunyi:
\begin{align*}
	\card{A} = \card{B} & \text{ jika dan hanya jika } A \approx B.
\intertext{Hal ini sangat mirip dengan prinsip yang kini disebut
\emph{Prinsip Hume}, yang menyatakan:}
	\fregenum{x} {F(x)} = \fregenum{x}{G(x)} & \text{ jika dan hanya jika } F \sim G
\end{align*}
di mana `$F \sim G$' menyingkat pernyataan bahwa objek-objek yang memenuhi
$F$ tepat sama banyaknya dengan objek-objek yang memenuhi $G$, yakni bahwa
kedua kumpulan itu dapat dipasangkan melalui suatu bijeksi, yakni:
\begin{align*}
	\exists R(&\forall v\forall y(Rvy \lif (Fv \land Gy)) \land {}\\
		&\forall v(Fv \lif \lexists![y][Rvy]) \land {}\\
		&\forall y(Gy \lif \lexists![v][Rvy]))
\end{align*}
Namun, terdapat perbedaan tipe antara Prinsip Hume dan Prinsip Cantor. Dalam
pernyataan Prinsip Cantor, variabel ``$A$'' dan ``$B$'' merupakan suku orde
pertama yang menyatakan \emph{himpunan}. Dalam pernyataan Prinsip Hume,
``$F$'', ``$G$'', dan ``$R$'' \emph{bukan} suku orde pertama; alih-alih,
semuanya berada dalam \emph{posisi predikat}. (Mungkin semuanya menyatakan
\emph{sifat}.) Jadi, kita dapat menguraikan Prinsip Hume dalam bahasa biasa
sebagai berikut: banyaknya objek yang memenuhi $F$ sama dengan banyaknya
objek yang memenuhi $G$ jika dan hanya jika kedua kumpulan objek itu dapat
dipasangkan secara bijektif. Prinsip ini disebut
\emph{Prinsip Hume} karena Hume pernah menulis:
\begin{quote}
  Apabila dua bilangan dipadankan sedemikian rupa sehingga bilangan yang satu
  selalu mempunyai suatu satuan yang bersesuaian dengan setiap satuan pada
  bilangan yang lain, kita menyatakan keduanya sama.
  \citep[Pt.III Bk.1 \S1]{Hume1740}
\end{quote}
Prinsip Hume diangkat ke kedudukan penting dalam matematika dan logika
kontemporer oleh \citet[\S63]{Frege1884}, yang mengutip bagian dari Hume ini
sebelum (pada dasarnya) menguraikan secara garis besar prinsip yang kita
sebut Prinsip Hume.

Anda patut memperhatikan kemiripan struktural antara Prinsip Hume dan Hukum
Dasar~V. Kita merumuskannya dalam \olref[story][blv]{sec} sebagai berikut:
\[
	\fregeext{x}{F(x)} = \fregeext{x}{G(x)} \text{ jika dan hanya jika } \lforall[x][(F(x) \liff G(x))].
\]
Pada titik ini, beberapa ulasan dan perbandingan mungkin membantu.

Terdapat dua cara memahami suatu prinsip seperti Prinsip Hume atau Hukum
Dasar~V: secara \emph{predikatif} atau secara \emph{impredikatif} (ingat
\olref[story][predicative]{sec}). Pada pembacaan impredikatif atas Hukum
Dasar~V, untuk setiap~$F$, objek $\fregeext{x}{F(x)}$ berada dalam domain
kuantifikasi yang kita gunakan untuk merumuskan Hukum Dasar~V itu sendiri.
Demikian pula, pada pembacaan impredikatif atas Prinsip Hume, untuk
setiap~$F$, objek $\fregenum{x}{F(x)}$ berada dalam domain kuantifikasi yang
kita gunakan untuk merumuskan Prinsip Hume. Sebaliknya, menurut pemahaman
\emph{predikatif}, objek $\fregeext{x}{F(x)}$ dan~$\fregenum{x}{F(x)}$ akan
menjadi entitas dari suatu domain yang \emph{berbeda}.

Jika kita membaca Hukum Dasar~V secara impredikatif, prinsip tersebut
menghasilkan ketidakkonsistenan melalui Komprehensi Naif (untuk rinciannya,
lihat \olref[story][blv]{sec}). Seperti Komprehensi Naif, prinsip tersebut
dapat dibuat konsisten dengan membacanya secara \emph{predikatif}. Namun,
kemungkinan besar prinsip itu tidak akan melakukan semua yang kita inginkan.

Akan tetapi, Prinsip Hume \emph{dapat} dibaca secara impredikatif dengan
konsisten. Dengan pembacaan tersebut, prinsip ini cukup kuat.

Sebagai ilustrasi, perhatikan predikat ``$x \neq x$'', yang jelas tidak
dipenuhi oleh apa pun. Prinsip Hume kini menghasilkan suatu objek
$\# x( x\neq x)$. Kita dapat memperlakukannya sebagai bilangan~$0$.
Sekarang, berdasarkan pemahaman \emph{impredikatif}---dan \emph{hanya}
berdasarkan pemahaman impredikatif---entitas $0$ ini berada dalam domain
kuantifikasi awal kita. Jadi, kita dapat secara bermakna menerapkan Prinsip
Hume dengan predikat ``$x = 0$'' untuk memperoleh suatu objek
$\#x (x = 0)$. Kita dapat memperlakukannya sebagai bilangan~$1$. Lebih jauh,
Prinsip Hume menyiratkan bahwa $0 \neq 1$, sebab tidak mungkin terdapat
bijeksi dari objek-objek yang tidak identik dengan dirinya sendiri ke
objek-objek yang identik dengan $0$ (tidak ada objek jenis pertama, tetapi
terdapat satu objek jenis kedua). Sekarang, dengan kembali bekerja secara
impredikatif, $1$~berada dalam domain kuantifikasi awal kita. Jadi, kita dapat
secara bermakna menerapkan Prinsip Hume dengan predikat
``$(x = 0 \lor x = 1)$'' untuk memperoleh suatu objek
$\#x(x = 0 \lor x = 1)$. Kita dapat memperlakukannya sebagai bilangan~$2$,
dan kita dapat menunjukkan bahwa $0\neq 2$, $1 \neq 2$, dan seterusnya.

Singkatnya, jika dipahami secara impredikatif, Prinsip Hume menyiratkan bahwa
terdapat \emph{tak hingga banyak objek}. Hal ini mendorong para
\emph{logisis neo-Fregean} untuk menjadikan Prinsip Hume sebagai landasan
aritmetika.

Namun, Frege \emph{sendiri} tidak menjadikan Prinsip Hume sebagai landasan
aritmetikanya. Sebaliknya, Frege membuktikan Prinsip Hume dari suatu definisi
eksplisit: $\fregenum{x}{F(x)}$ didefinisikan sebagai ekstensi konsep
$F \sim \Phi$. Dalam istilah modern, kita mungkin mencoba menyatakannya
sebagai $\fregenum{x}{F(x)} = \Setabs{G}{F \sim G}$; tetapi langkah ini akan
membawa kita kembali ke masalah-masalah Komprehensi Naif.

\end{document}
